\section[Detailed computations for the S3 obstruction]{Detailed computations for the \texorpdfstring{$S_3$}{S3} obstruction}
\label{app:s3-model-computations}

This appendix contains the matrix calculations and coefficient extractions used in \Cref{sec:shadows}. The main text keeps only the structural obstruction theorem; the computational certificates are recorded here to keep the quotient-shadow package from obscuring the Adams--M\"obius class-Mertens spine.


\begin{proposition}[Quadratic Adams obstruction]
  \label{thm:quadratic-adams-obstruction}
  \label{prop:quadratic-adams-obstruction}
  Let \(\chi\) be a non-trivial quadratic character of \(G\). Then, for
  \(|z|<\lambda^{-1}\),
  \[
    \Pi_\chi(z)
    =
    \sum_{\substack{m\ge 1\\ m\ \mathrm{odd}}}
      \frac{\mu(m)}{m}\,\mathcal L_\chi(z^m)
    +
    \sum_{\substack{m\ge 1\\ m\ \mathrm{even}}}
      \frac{\mu(m)}{m}\,\mathcal L_{\mathbf 1}(z^m).
  \]
  In particular, if \(\mathcal L_\chi\equiv0\), then
  \[
    \Pi_\chi(z)
    =
    -\frac{1}{2}
      \sum_{\substack{r\ge 1\\ r\ \mathrm{odd}}}
        \frac{\mu(r)}{r}\,
        \mathcal L_{\mathbf 1}(z^{2r}),
  \]
  and hence \(\Pi_\chi\) need not vanish.
\end{proposition}

\begin{proof}
  Since \(\chi\) is quadratic, \(\chi^m=\chi\) for \(m\) odd and
  \(\chi^m=\mathbf1\) for \(m\) even. The Adams--M\"obius inversion
  formula \eqref{eq:class-adams-mobius-inversion} gives the first
  identity. If \(\mathcal L_\chi\equiv0\), only even indices remain;
  writing \(m=2r\) and using \(\mu(2r)=-\mu(r)\) for odd squarefree
  \(r\) gives the second identity.
\end{proof}

Consider the full two-shift
\[
  A=
  \begin{pmatrix}
    1 & 1 \\
    1 & 1
  \end{pmatrix},
  \qquad
  \lambda=\rho(A)=2,
\]
and label the four allowed edges by
\begin{equation}\label{eq:s3-example-edge-labels}
  \begin{aligned}
    \tau(0\to 0)&=(23), & \tau(0\to 1)&=e, \\
    \tau(1\to 0)&=(12), & \tau(1\to 1)&=(123).
  \end{aligned}
\end{equation}
Let \(T\) be the transposition class and \(K\) the \(3\)-cycle class in
\(S_3\). The character table is
\begin{equation}\label{eq:s3-character-table}
  \begin{array}{c|ccc}
    & e & T & K \\ \hline
    \mathbf 1 & 1 & 1 & 1 \\
    \varepsilon & 1 & -1 & 1 \\
    \rho & 2 & 0 & -1
  \end{array}
\end{equation}
where \(\varepsilon\) is the sign character and \(\rho\) the standard
two-dimensional representation. In this \(S_3\) example only, \(\rho\)
also denotes the standard representation; expressions such as
\(\rho(M)\) still denote spectral radius.

In the basis
\[
  v_1=e_1-e_2,\qquad v_2=e_1+e_2-2e_3,
\]
the standard representation matrices of the generators appearing in
\eqref{eq:s3-example-edge-labels} are
\begin{equation}\label{eq:s3-standard-representation-matrices}
\begin{aligned}
  \rho((12))
    &=
    \begin{pmatrix}
      -1 & 0 \\
      0 & 1
    \end{pmatrix}, \\
  \rho((23))
    &=
    \begin{pmatrix}
      \frac12 & \frac32 \\
      \frac12 & -\frac12
    \end{pmatrix}, \\
  \rho((123))
    &=
    \begin{pmatrix}
      -\frac12 & -\frac32 \\
      \frac12 & -\frac12
    \end{pmatrix}.
\end{aligned}
\end{equation}
These matrices are written in a convenient equivalent real basis, not
necessarily a unitary one; the character identities, conjugation
relations, and spectral-radius statements used below are invariant under
equivalence of representations.

\begin{proposition}[Twisted channels for the \(S_3\)-example]
  \label{prop:s3-example-channels}
  One has
  \begin{equation}\label{eq:s3-example-aeps}
    A_\varepsilon=
    \begin{pmatrix}
      -1 & 1 \\
      -1 & 1
    \end{pmatrix},
    \qquad
    A_\varepsilon^2=0,
    \qquad
    \det(I-zA_\varepsilon)=1.
  \end{equation}
  Moreover,
  \begin{equation}\label{eq:s3-example-arho}
    A_\rho=
    \begin{pmatrix}
      \frac12 & \frac32 & 1 & 0 \\
      \frac12 & -\frac12 & 0 & 1 \\
      -1 & 0 & -\frac12 & -\frac32 \\
      0 & 1 & \frac12 & -\frac12
    \end{pmatrix},
  \end{equation}
  and
  \begin{equation}\label{eq:s3-example-arho-charpoly}
    \chi_{A_\rho}(t)=t^3(t+1),
    \qquad
    \det(I-zA_\rho)=1+z.
  \end{equation}
  Consequently,
  \begin{equation}\label{eq:s3-example-logarithms}
    \mathcal L_\varepsilon(z)=0,
    \qquad
    \mathcal L_\rho(z)=\log\det(I-zA_\rho)^{-1}=-\log(1+z).
  \end{equation}
  If \(a_{n,C}\) denotes the number of period-\(n\) points whose
  \(S_3\)-label lies in \(C\), then
  \begin{equation}\label{eq:s3-example-periodic-classes}
    a_{n,e}=\frac{2^n+2(-1)^n}{6},
    \qquad
    a_{n,T}=2^{n-1},
    \qquad
    a_{n,K}=\frac{2^n-(-1)^n}{3}.
  \end{equation}
  In particular, the transposition class is exactly equidistributed at
  every period length:
  \begin{equation}\label{eq:s3-example-periodic-equidistribution}
    a_{n,T}=\frac{|T|}{|S_3|}\,2^n
    \qquad \text{for every } n\ge 1.
  \end{equation}
  Moreover,
  \begin{equation}\label{eq:s3-example-eta-check}
    \rho(A_\varepsilon)=0,
    \qquad
    \rho(A_\rho)=1,
    \qquad
    \eta=1<2=\lambda.
  \end{equation}
  In particular, the hypothesis of
  \Cref{thm:class-mertens-explicit} holds for this example.
\end{proposition}

\begin{proof}
  The sign matrix is obtained by substituting
  \(\varepsilon((23))=-1\), \(\varepsilon(e)=1\),
  \(\varepsilon((12))=-1\), and \(\varepsilon((123))=1\) into the
  \(2\times2\) block matrix of edge labels. This gives
  \[
    A_\varepsilon
    =
    \begin{pmatrix}
      -1 & 1 \\
      -1 & 1
    \end{pmatrix},
    \qquad
    A_\varepsilon^2=0,
  \]
  and therefore \(\det(I-zA_\varepsilon)=1\).

  For the standard channel, \eqref{eq:s3-standard-representation-matrices}
  yields the block form
  \[
    A_\rho
    =
    \begin{pmatrix}
      \rho((23)) & I_2 \\
      \rho((12)) & \rho((123))
    \end{pmatrix},
  \]
  which is exactly \eqref{eq:s3-example-arho}. To compute its
  characteristic polynomial, write
  \[
    tI_4-A_\rho
    =
    \begin{pmatrix}
      tI_2-\rho((23)) & -I_2 \\
      -\rho((12)) & tI_2-\rho((123))
    \end{pmatrix}.
  \]
  Since the upper-right block \(-I_2\) commutes with the lower-right
  block, the block determinant identity gives
  \[
    \det(tI_4-A_\rho)
    =
    \det\!\bigl((tI_2-\rho((123)))(tI_2-\rho((23)))-\rho((12))\bigr).
  \]
  A direct multiplication using
  \eqref{eq:s3-standard-representation-matrices} shows that
  \[
    (tI_2-\rho((123)))(tI_2-\rho((23)))-\rho((12))
    =
    \begin{pmatrix}
      t^2 & 0 \\
      -t & t(t+1)
    \end{pmatrix}.
  \]
  Hence
  \[
    \chi_{A_\rho}(t)=\det(tI_4-A_\rho)=t^3(t+1),
  \]
  so the eigenvalues of \(A_\rho\) are \(0,0,0,-1\). This gives
  \(\rho(A_\rho)=1\), \(\det(I-zA_\rho)=1+z\), and
  \(\Tr(A_\rho^n)=(-1)^n\) for every \(n\ge 1\).

  Since \(\Tr(A^n)=2^n\), \(\Tr(A_\varepsilon^n)=0\), and
  \(\Tr(A_\rho^n)=(-1)^n\), the logarithmic channels are
  \[
  \begin{aligned}
    \mathcal L_\varepsilon(z)
      &=\log\det(I-zA_\varepsilon)^{-1}=0,\\
    \mathcal L_\rho(z)
      &=\log\det(I-zA_\rho)^{-1}=-\log(1+z).
  \end{aligned}
  \]
  The class-counting formula from
  \Cref{prop:class-indicator-expansion} gives
  \[
    \begin{aligned}
      a_{n,C}
      &=
      \frac{|C|}{|S_3|}
        \Bigl(
          \Tr(A^n)
          +\chi_\varepsilon^\ast(C)\Tr(A_\varepsilon^n)
          +\chi_\rho^\ast(C)\Tr(A_\rho^n)
        \Bigr) \\
      &=
      \frac{|C|}{|S_3|}
        \Bigl(2^n+\chi_\rho^\ast(C)(-1)^n\Bigr),
    \end{aligned}
  \]
  from which \eqref{eq:s3-example-periodic-classes} follows.
\end{proof}

\begin{proposition}[Primitive channels for the \(S_3\)-example]
  \label{prop:s3-example-primitive-channels}
  The sign channel satisfies
  \begin{equation}\label{eq:s3-example-pi-eps}
    \Pi_\varepsilon(z)
    =
    -\frac12
      \sum_{\substack{r\ge 1\\ r\ \mathrm{odd}}}
        \frac{\mu(r)}{r}\log(1-2z^{2r})^{-1},
  \end{equation}
  with the Perron-point certificate recorded in
  \Cref{lem:s3-example-perron-certificates}.

  For the standard representation,
  \begin{equation}\label{eq:s3-example-adams-rho}
    \psi^m\rho=
    \begin{cases}
      \rho, & m\equiv \pm 1 \pmod 6, \\
      \mathbf 1-\varepsilon+\rho, & m\equiv \pm 2 \pmod 6, \\
      \mathbf 1+\varepsilon, & m\equiv 3 \pmod 6, \\
      2\mathbf 1, & 6\mid m,
    \end{cases}
  \end{equation}
  and therefore
  \begin{equation}\label{eq:s3-example-pi-rho}
    \Pi_\rho(z)
    =
    \sum_{m\ge 1}\frac{\mu(m)}{m}
      \Bigl(
        \alpha(m)\log(1-2z^m)^{-1}
        -\beta(m)\log(1+z^m)
      \Bigr),
  \end{equation}
  where
  \begin{equation}\label{eq:s3-example-alpha-beta}
    (\alpha(m),\beta(m))=
    \begin{cases}
      (0,1), & m\equiv \pm 1 \pmod 6, \\
      (1,1), & m\equiv \pm 2 \pmod 6, \\
      (1,0), & m\equiv 3 \pmod 6, \\
      (2,0), & 6\mid m,
    \end{cases}
  \end{equation}
  with the Perron-point certificate again recorded in
  \Cref{lem:s3-example-perron-certificates}.
\end{proposition}

\begin{proof}
  Since \(\varepsilon\) is quadratic and
  \(\mathcal L_\varepsilon(z)=0\), \Cref{prop:quadratic-adams-obstruction}
  gives \eqref{eq:s3-example-pi-eps}. The series converges absolutely
  at \(z=1/2\), and the certified non-zero interval is established in
  \Cref{lem:s3-example-perron-certificates}.

  For \(\rho\), the Adams powers are computed from the character table
  \eqref{eq:s3-character-table} together with the power maps on the
  three conjugacy classes. In this model those maps are immediate: a
  transposition has \(m\)-th power in \(T\) for \(m\) odd and in
  \(\{e\}\) for \(m\) even, while a \(3\)-cycle has \(m\)-th power in
  \(K\) when \(3\nmid m\) and in \(\{e\}\) when \(3\mid m\). Therefore
  the class values of \(\psi^m\rho\) are
  \[
    \begin{array}{c|ccc|c}
      m \bmod 6
      & (\psi^m\chi_\rho)(e)
      & (\psi^m\chi_\rho)(T)
      & (\psi^m\chi_\rho)(K)
      & \psi^m\rho \\ \hline
      \pm 1 & 2 & 0 & -1 & \rho \\
      \pm 2 & 2 & 2 & -1 & \mathbf 1-\varepsilon+\rho \\
      3 & 2 & 0 & 2 & \mathbf 1+\varepsilon \\
      0 & 2 & 2 & 2 & 2\mathbf 1
    \end{array}
  \]
  and decomposing these four class functions in the basis
  \(\{\mathbf 1,\varepsilon,\rho\}\) gives
  \eqref{eq:s3-example-adams-rho}. Substituting the decomposition into
  \Cref{cor:primitive-peter-weyl-determinant}, and using
  \eqref{eq:s3-example-logarithms}, yields
  \eqref{eq:s3-example-pi-rho} and \eqref{eq:s3-example-alpha-beta}.
  The resulting series again converges absolutely at \(z=1/2\), and the
  certified non-zero interval is established in
  \Cref{lem:s3-example-perron-certificates}.
\end{proof}

\begin{proposition}[Naive Artin formula obstruction at the Euler-product
  level]
  \label{prop:s3-naive-artin-obstruction}
  In the \(S_3\)-example,
  \begin{equation}\label{eq:s3-f-level-windows}
    -\frac{381}{1000}<F_\varepsilon(1/2)<-\frac{380}{1000},
    \qquad
    -\frac{923}{1000}<F_\rho(1/2)<-\frac{922}{1000}.
  \end{equation}
  At the same point the ordinary Artin logarithm in the sign channel is
  identically zero:
  \begin{equation}\label{eq:s3-artin-zero-f-nonzero}
    \mathcal L_\varepsilon(1/2)=0,
    \qquad
    F_\varepsilon(1/2)\neq0.
  \end{equation}
  Hence Frobenius-class Euler-product constants cannot in general be
  obtained by replacing the fixed-label transforms in
  \Cref{thm:class-euler-product-mertens} with ordinary Artin logarithms.
\end{proposition}

\begin{proof}
  The equality \(\mathcal L_\varepsilon(z)\equiv0\) is
  \eqref{eq:s3-example-logarithms}. The certified intervals for
  \(F_\varepsilon(1/2)\) and \(F_\rho(1/2)\) are proved by exact rational
  interval arithmetic in \Cref{lem:appendix-s3-f-level-certificates}; in
  particular both intervals exclude zero. Since
  \Cref{thm:class-euler-product-mertens} uses
  \(F_\varrho(\lambda^{-1})\), whereas
  \Cref{thm:fixed-label-euler-transform} identifies
  \(\mathcal L_\varrho(z)=\sum_{r\ge1}\Pi_{\psi^r\varrho}(z^r)/r\), the
  certified sign-channel inequality gives a concrete obstruction to the
  naive Artin substitution.
\end{proof}

\begin{lemma}[Certified \(S_3\) non-vanishing]
  \label{lem:s3-example-perron-certificates}
  For the \(S_3\)-extension above, the Perron-point primitive channels
  satisfy the certified inequalities
  \begin{equation}\label{eq:s3-example-certified-intervals}
    -\frac{342}{1000}
    <\Pi_\varepsilon(1/2)<-\frac{341}{1000},
    \qquad
    -\frac{719}{1000}
    <\Pi_\rho(1/2)<-\frac{718}{1000}.
  \end{equation}
  Hence both primitive channels are bounded away from zero, and the
  higher-dimensional term in the non-abelian defect from
  \Cref{thm:abelian-shadow-defect} is genuinely present in this
  example.
\end{lemma}

\begin{proof}
  Appendix~\ref{app:s3-log-certificates} gives a self-contained
  rational interval calculation for the finite truncations and tails of
  the two logarithmic series. Its enclosures imply exactly the coarse
  windows in \eqref{eq:s3-example-certified-intervals}. Since neither
  interval contains \(0\), both primitive channels are non-zero, and the
  concluding statement follows from
  \Cref{thm:abelian-shadow-defect}.
\end{proof}

\begin{corollary}[Primitive profile and deviations for the
  \(S_3\)-example]
  \label{cor:s3-example-profile}
  For the \(S_3\)-extension defined above, one has
  \begin{equation}\label{eq:s3-example-primitive-sign-coefficients}
    \pi_{n,\varepsilon}(A,\tau)
    =
    \frac{1}{n}
      \sum_{\substack{m\mid n\\ 2\mid m}}
        \mu(m)\,2^{n/m},
  \end{equation}
  and
  \begin{equation}\label{eq:s3-example-primitive-rho-coefficients}
    \pi_{n,\rho}(A,\tau)
    =
    \frac{1}{n}
      \sum_{m\mid n}
        \mu(m)\Bigl(
          \alpha(m)2^{n/m}+\beta(m)(-1)^{n/m}
        \Bigr).
  \end{equation}
  Consequently,
  \begin{equation}\label{eq:s3-example-primitive-classes}
    \begin{aligned}
      p_{n,e}
      &=\frac{p_n(A)+\pi_{n,\varepsilon}(A,\tau)+2\pi_{n,\rho}(A,\tau)}{6}, \\
      p_{n,T}
      &=\frac{p_n(A)-\pi_{n,\varepsilon}(A,\tau)}{2}, \\
      p_{n,K}
      &=\frac{p_n(A)+\pi_{n,\varepsilon}(A,\tau)-\pi_{n,\rho}(A,\tau)}{3}.
    \end{aligned}
  \end{equation}
  where
  \begin{equation}\label{eq:s3-example-full-two-shift-primitive}
    p_n(A)=\frac{1}{n}\sum_{d\mid n}\mu(d)\,2^{n/d}.
  \end{equation}
  The first primitive class triples are
  \begin{equation}\label{eq:s3-example-primitive-first-terms}
    (p_{n,e},p_{n,T},p_{n,K})
    =(0,1,1),(0,1,0),(0,1,1),(0,2,1),(1,3,2),(1,5,3),\ldots
  \end{equation}
  Finally,
  \begin{equation}\label{eq:s3-example-delta-values}
    \begin{aligned}
      \Delta_e(A,\tau;S_3)
      &=\frac16\Pi_\varepsilon(1/2)+\frac13\Pi_\rho(1/2), \\
      \Delta_T(A,\tau;S_3)
      &=-\frac12\Pi_\varepsilon(1/2), \\
      \Delta_K(A,\tau;S_3)
      &=\frac13\Pi_\varepsilon(1/2)-\frac13\Pi_\rho(1/2).
    \end{aligned}
  \end{equation}
  and the Perron-point windows imply
  \begin{equation}\label{eq:s3-example-delta-certified-windows}
    -\frac{89}{300}
    <\Delta_e(A,\tau;S_3)<-\frac{1777}{6000},
    \qquad
    \frac{341}{2000}
    <\Delta_T(A,\tau;S_3)<\frac{171}{1000},
  \end{equation}
  \begin{equation}\label{eq:s3-example-delta-k-certified-window}
    \frac{47}{375}
    <\Delta_K(A,\tau;S_3)<\frac{63}{500}.
  \end{equation}
  Equivalently, the class-constant deviations split as
  \begin{equation}\label{eq:s3-example-abelian-shadow}
    \Delta^{\mathrm{ab}}
    =
    \Bigl(
      \frac16\Pi_\varepsilon(1/2),
      -\frac12\Pi_\varepsilon(1/2),
      \frac13\Pi_\varepsilon(1/2)
    \Bigr),
  \end{equation}
  \begin{equation}\label{eq:s3-example-nonabelian-defect}
    \Delta^{\mathrm{na}}
    =
    \Bigl(
      \frac13\Pi_\rho(1/2),
      0,
      -\frac13\Pi_\rho(1/2)
    \Bigr).
  \end{equation}
  In particular,
  \begin{equation}\label{eq:s3-example-energy-ratio}
    \frac{\norm{\mathcal B_{A,\tau;G}^{\mathrm{na}}}_G^2}
         {\norm{\mathcal B_{A,\tau;G}^{\mathrm{ab}}}_G^2}
    =
    \frac{\abs{\Pi_\rho(1/2)}^2}
         {\abs{\Pi_\varepsilon(1/2)}^2}
    >4,
  \end{equation}
  so that, in the normalized \(B\)-profile \(L^2\)-norm of
  \Cref{thm:abelian-shadow-defect}, the genuine non-abelian primitive
  bias is certified to be more than four times stronger than the
  abelian shadow.
\end{corollary}

\begin{proof}
  Expanding \(\log(1-2z^{2r})^{-1}=\sum_{k\ge 1}2^k z^{2rk}/k\) in
  \eqref{eq:s3-example-pi-eps} and reading off the coefficient of
  \(z^n\) gives \eqref{eq:s3-example-primitive-sign-coefficients}.
  Likewise, coefficient extraction in \eqref{eq:s3-example-pi-rho}
  yields \eqref{eq:s3-example-primitive-rho-coefficients}. To obtain the
  three class formulas explicitly, specialise
  \Cref{prop:class-primitive-decomposition} to the character table
  \eqref{eq:s3-character-table}. Since the class sizes are
  \(|e|=1\), \(|T|=3\), and \(|K|=2\), and since
  \(\pi_{n,\mathbf1}(A,\tau)=p_n(A)\), character orthogonality gives
  \[
    \begin{aligned}
    p_{n,e}
      &=\frac{1}{6}\bigl(p_n(A)+\pi_{n,\varepsilon}(A,\tau)
        +2\pi_{n,\rho}(A,\tau)\bigr),\\
    p_{n,T}
      &=\frac{3}{6}\bigl(p_n(A)-\pi_{n,\varepsilon}(A,\tau)
        +0\cdot\pi_{n,\rho}(A,\tau)\bigr),\\
    p_{n,K}
      &=\frac{2}{6}\bigl(p_n(A)+\pi_{n,\varepsilon}(A,\tau)
        -\pi_{n,\rho}(A,\tau)\bigr),
    \end{aligned}
  \]
  which is \eqref{eq:s3-example-primitive-classes}. The scalar primitive
  count \eqref{eq:s3-example-full-two-shift-primitive} is the usual
  M\"obius formula for the full two-shift. Evaluating these formulas for
  \(n=1,\dots,6\) yields
  \eqref{eq:s3-example-primitive-first-terms}.

  The constant formulas are obtained from
  \Cref{prop:class-mertens-deviations} and the character table
  \eqref{eq:s3-character-table}, namely
  \[
    \begin{aligned}
      \Delta_e(A,\tau;S_3)
      &=\frac{1}{6}\Pi_\varepsilon(1/2)+\frac{1}{3}\Pi_\rho(1/2), \\
      \Delta_T(A,\tau;S_3)
      &=-\frac{1}{2}\Pi_\varepsilon(1/2), \\
      \Delta_K(A,\tau;S_3)
      &=\frac{1}{3}\Pi_\varepsilon(1/2)-\frac{1}{3}\Pi_\rho(1/2).
    \end{aligned}
  \]
  Combining the certified intervals in
  \Cref{lem:s3-example-perron-certificates} with these formulas gives
  \eqref{eq:s3-example-delta-certified-windows} and
  \eqref{eq:s3-example-delta-k-certified-window}. The signs and
  non-vanishing of these three deviations follow immediately.
  The decomposition into
  \eqref{eq:s3-example-abelian-shadow} and
  \eqref{eq:s3-example-nonabelian-defect} is exactly the splitting from
  \Cref{thm:abelian-shadow-defect}. Finally,
  \eqref{eq:s3-example-energy-ratio} is the corresponding norm identity,
  and the bound \(>4\) follows from
  \eqref{eq:s3-example-certified-intervals}, since
  \((718/342)^2>4\). The following table is therefore only a compact
  summary of the identities just derived and of the certified windows used
  in the last step.
\end{proof}

\begin{center}
\footnotesize
\begin{tabular}{@{}p{0.16\textwidth}p{0.32\textwidth}p{0.36\textwidth}@{}}
\toprule
Channel / observable & Periodic-level behaviour & Certified primitive information \\
\midrule
\(\varepsilon\) &
\(\mathcal L_\varepsilon(z)=0\) &
\(-\frac{342}{1000}<\Pi_\varepsilon(1/2)<-\frac{341}{1000}\) \\
\(\rho\) &
\(\mathcal L_\rho(z)=-\log(1+z)\) &
\(-\frac{719}{1000}<\Pi_\rho(1/2)<-\frac{718}{1000}\) \\
\(T\) &
\(a_{n,T}=2^{n-1}\) for all \(n\) &
\(\frac{341}{2000}<\Delta_T(A,\tau;S_3)<\frac{171}{1000}\) \\
\bottomrule
\end{tabular}
\end{center}
